% !PRINTED 657
\weberpage{657}
%<<CONT

und hierin liegt der Beweis des folgenden Hauptsatzes, den man
erhält, wenn man $t = \tau_1$ setzt.

6. \gesperrt{Wenn $e_1 = 1$ ist, so ist $F'(\tau)$ nicht durch $\pi_1$
teilbar. Sonst ist $F'(\tau)$ durch $\pi_1^{e_1-1}$, aber nicht durch
$\pi_1^{e_1}$ teilbar. Das Gleiche gilt von den übrigen Primfaktoren
von $(z-c)$ und folglich ist $F'(\tau)$ teilerfremd zu allen
Primfaktoren $z-c$ in $Z$, die nicht durch das Quadrat eines
Primfaktors in $\Omega$ teilbar sind.}

Es gibt also nur eine endliche Anzahl von Primfaktoren $(z-c)$, die
durch eine höhere als die erste Potenz eines Primfaktors in $\Omega$
teilbar sind, und $F'(\tau)$ ist, in Primfaktoren zerlegt, das
Produkt aller Faktoren $\pi^{e-1}$, wenn $\pi^e$ die höchste in
$N(\pi) = z - c$ aufgehende Potenz von $\pi$ ist.

Das Funktional $F'(\tau)$ wird das \gesperrt{Verzweigungsfunktional} des
Körpers $\Omega$ genannt.

Für die absolute Norm des Verzweigungsfunktionals erhalten wir:

\weq{24}{N_a F'(\tau) = \Pi (z-c)^{e_1-1+e_2-1+\cdots},}

worin sich das Produkt $\Pi$ auf alle Linearfaktoren $(z-c)$
erstreckt, die durch eine höhere als die erste Potenz eines
Primfaktors teilbar sind.

Wir können noch zeigen, daß diese absolute Norm nichts anderes ist,
als die \gesperrt{Körperdiskriminante} $D$.

Um dies zu beweisen, setzen wir:

\weq{25}{\tau^k = u_{k,1}\omega_1 + u_{k,2}\omega_2 + \cdots + u_{k,n}\omega_n,}
% [sic] printed u_{k,1}\omega_1 + ... + u_{k,n}\omega_n (and u_{k,1}, ..., u_{k,n} in the next sentence) with the power index k first, inconsistent with (26) and the congruence below, which put the power index second (u_{1,0}, ..., u_{n,n-1}; u_{r,0}y_0 + ...)

worin die $u_{k,1}, u_{k,2}, \ldots, u_{k,n}$ homogene Funktionen
$k$ten Grades von $t_1, t_2, \ldots, t_n$ sind, und die Determinante

\weq{26}{U =
\begin{vmatrix}
u_{1,0}, & u_{2,0}, & \ldots, & u_{n,0} \\
u_{1,1}, & u_{2,1}, & \ldots, & u_{n,1} \\
\cdot & \cdot & \cdot & \cdot \\
u_{1,n-1}, & u_{2,n-1}, & \ldots, & u_{n,n-1}
\end{vmatrix}}

ist eine Einheit.

Denn wäre $U$ keine Einheit, so hätte sie wenigstens einen
Linearfaktor $z-c$, und man könnte die ganzen rationalen holomorphen
Funktionale $y_0, y_1, \ldots, y_{n-1}$ so bestimmen, daß

\[ u_{r,0}y_0 + u_{r,1}y_1 + \cdots + u_{r,n-1}y_{n-1} \equiv 0 \quad (\mathrm{mod}\ z-c) \]

wäre, ohne daß alle $y_s$ durch $z-c$ teilbar sind.
